\hypertarget{network-egress}{%
\chapter{6. Network egress}\label{network-egress}}

The network is how a workload reaches the world, and how the world
reaches back. Give a workload the host's network untouched and it has
all of it: it can bind ports, open a connection to anywhere on the
internet, sweep the local subnet, and reach the private services the
host runs only for itself. A workload that has been turned reaches for
exactly these: a channel out to a controller, a route for the user's
data to leave by, a path sideways into the home network and the services
on it (\texttt{T1.6}).

The network is also the one resource a working agent cannot simply be
denied. It has to reach a package registry, an API, a source it was told
to build from. So the network is the first place both limbs are needed
at once: absence, to take away the network the workload should never
have, and interposition, to meter the narrow slice it needs.

\hypertarget{absence-then-interposition}{%
\section{6.1 Absence, then
interposition}\label{absence-then-interposition}}

Absence comes first, and it does most of the work. In the workload's
world the host's network is not there: no interfaces, no routes, no
inherited connections, none of the private services the host runs only
for itself. A workload mapping its surroundings for an administrative
port, or for a service the host meant only for its own use, finds
nothing there to map (\texttt{T1.6}); the lateral-movement surface is
not defended, it is absent. Most of what a network boundary has to stop
never has to be stopped here, because there is nothing to connect to and
no route by which to try. The forbidden connection cannot be formed.

What absence cannot do is give the workload its registry back. The
workload needs a route out, and a route out is exactly what was removed.
So the egress that remains, the connection to the outside the workload
is allowed to make, is carried by the second limb. The empty network has
no path to the internet; what it has is the monitor. Each outbound
connection is a request the workload makes and the monitor answers
(\texttt{interpose-as-transaction}): the workload names a destination,
the monitor checks it against the policy, and only if it passes does the
monitor open the connection and hand it back. A connection the policy
does not allow is one the monitor never opens, and the workload cannot
reach past it to dial for itself, because it has no route except through
the request.

The workload names its destination; it never resolves it. It has no
resolver to reach and no address of its own to offer: it asks for a
name, and the monitor resolves that name and pins the address it
resolved to before dialing. This is what closes off rebinding: a forged
answer cannot swing an approved name onto a forbidden address, because
the address the monitor dials is the address the monitor checked, and
the workload was never given the chance to supply one. The decision
falls on the destination that will be connected to, not the one that was
asked for.

And the monitor decides without carrying. Once a connection is approved
the bytes flow directly between the workload and the destination, the
monitor out of the path entirely (\texttt{control-not-data-plane}). It
is the point every connection is decided at, and no part of what any
connection moves. That is what keeps it small enough to trust: a thing
that ruled on a million bytes would have to be trusted with a million
bytes; a thing that rules only on whether a connection may open has
almost nothing in it to get wrong.

\hypertarget{the-contract}{%
\section{6.2 The contract}\label{the-contract}}

The contract is a spectrum, anchored at the safe end. A workload is
given no network at all until the policy asks for one, and asked for
conservatively unless the operator widens it: deny-by-default applied to
reach, the same fail-safe defaults the filesystem and execution
boundaries rest on {[}SaltzerSchroeder{]}.

At the closed end is the empty network with nothing added back: no
connection, no bind, no surface. It is where untrusted analysis and
post-install scripts belong, things that have no business touching the
network at all. A step out from there grants the isolated network a
named set of destinations and nothing else: the workload reaches the
registry and the API it was told it could reach, by asking for them, and
reaches nothing else, because nothing else has been named. Wider still
opens egress to the open internet, for the workload that must fetch from
sources no allowlist can predict; the network is still its own, the
connections are still brokered, the floor still holds. And at the open
end is the host's network stack, shared directly: no separate network,
no broker, the host's local services reachable once more. This last is
the escape hatch, and it reinstates the very surface the chapter began
by removing (\texttt{T1.6}), so the design will not compile it without
an explicit reason written down beside it
(\texttt{footgun-warn-dont-forbid}).

Beneath every posture but the empty one runs a floor no policy can lift.
A handful of destinations are denied in every mode, the widest included:
the address a cloud platform answers its own metadata on and the
link-local ranges, the standing targets of request-forgery and
infrastructure reconnaissance. Even the open-internet posture cannot
reach them, because the denial is not a policy choice but a property of
the boundary. The broker holds one more line of the same kind: it will
not become an open relay. A workload cannot hand it an arbitrary
internal address and have it dialed, not a sibling's network, not the
host's own services. Reaching a host service is its own named grant,
deliberate and visible, never something egress backs into.

\hypertarget{what-the-boundary-does-not-claim}{%
\section{6.3 What the boundary does not
claim}\label{what-the-boundary-does-not-claim}}

The boundary governs where a workload connects. It does not govern what
the workload says once connected, and the line between those two is not
a corner the design cut: it is the edge of what a boundary of this kind
can enforce at all.

That a workload never opens a connection the policy forbids is a safety
property: the forbidden connection is a thing that, once it happens,
cannot be called back, and the monitor's whole purpose is to see that it
does not happen. Complete mediation of the channels a workload may open
is enforceable, and the monitor enforces it. What travels over a channel
that was allowed to open is a different kind of thing. It is information
flow, and information flow is not a safety property: no monitor watching
one execution at a time can enforce it {[}Schneider{]}. So the monitor
does not look. There is no payload inspection and no interception of
encrypted streams; the bytes over an approved connection are opaque to
it by design.

This draws the residual exactly. When an operator allows a workload to
reach an external API, what the workload chooses to send there is a
data-governance question, not a confinement one (\texttt{T1.8}). The
boundary mediates the channel completely and makes no claim about the
payload. The claim it does not make is one no boundary of this shape
could honour, which is why stating it plainly is the honest thing rather
than a confession. The host-mode residual is of the same character: an
operator who shares the host's stack has reopened the lateral surface on
purpose (\texttt{T1.6}), and the boundary records the choice rather than
overriding it.

Across all of it the two limbs divide the surface the way the third
chapter said they would. Absence takes the network the workload should
never have had and leaves nothing to mediate; interposition meters the
one slice that could not be removed, and meters it at the point of
decision without ever touching the flow. Together they enforce a clean
safety property (no forbidden connection occurs) and they stop exactly
where safety stops, at the payload, which was never theirs to police.
